\clearpage
\hypertarget{native-264-pure13}{}\label{native-264-pure13}
\Article{13}
Da wir nun gezeigt haben, wie sich ein Produkt aus zwei Faktoren entwickeln lässt, wird die Entwicklung eines Produkts aus mehr Faktoren keine Schwierigkeit bereiten. Ist das Produkt $[r,R^\mu]\times[r,R^\nu]$ auf die Form $A[1,R^{\mu+\nu}]+B[r,R^{\mu+\nu}]$ gebracht, so wird beim Hinzufügen eines dritten Faktors $[r,R^\pi]$ das Produkt offenbar $=C[1,R^{\mu+\nu+\pi}]+D[r,R^{\mu+\nu+\pi}]$, wenn wir
\[
[r,R^{\mu+\nu}][r,R^\pi]=c[1,R^{\mu+\nu+\pi}]+d[r,R^{\mu+\nu+\pi}]
\]
und
\[
\begin{aligned}
C&=Bc\\
D&=Bd+A\{1+R^{\mu+\nu}+R^{2\mu+2\nu}+\text{usw.}+R^{(n-2)(\mu+\nu)}\}
\end{aligned}
\]
setzen. Somit lässt sich $[r,R]^\lambda$ leicht auf die Form $A[1,R^\lambda]+B[r,R^\lambda]$ bringen.

Als Beispiel wollen wir die Potenzen der Funktion $[r,R]$ für $n=11$, $\beta=5$ entwickeln und dabei $g=2$ setzen. Den Zahlen entsprechen dann die folgenden Indizes:
\[
\begin{array}{r@{\quad}r@{\pdot}r@{\pdot}r@{\pdot}r@{\pdot}r@{\pdot}r@{\pdot}r@{\pdot}r@{\pdot}r@{\pdot}r}
\text{Zahlen}&1&2&3&4&5&6&7&8&9&10\\
\text{Indizes}&0&1&8&2&4&9&7&3&6&5
\end{array}
\]
Zur Entwicklung des Quadrats $[r,R]^2$ nach Formel I aus Artikel~11 haben wir also
\[
\mu=1,\quad\nu=1
\]
\[
\begin{array}{l@{\quad}r@{\pdot}r@{\pdot}r@{\pdot}r@{\pdot}r@{\pdot}r@{\pdot}r@{\pdot}r@{\pdot}r}
\text{Werte von }x\ \ldots&1&2&3&4&5&6&7&8&9\\
\operatorname{ind}x\ \ldots&0&1&8&2&4&9&7&3&6\\
2\operatorname{ind}(x+1)\ \ldots&2&16&4&8&18&14&6&12&10\\
\multicolumn{10}{l}{\text{Kleinste Reste von }\operatorname{ind}x-2\operatorname{ind}(x+1)}\\
\quad\text{zum Modul }5\ \ldots&3&0&4&4&1&0&1&1&1
\end{array}
\]
woraus wir
\SourceLeaf{265}
\[
\Omega=[r,R^2]\times\{2+4R+R^3+2R^4\}
\]
und
\[
1^0.\qquad[r,R]^2=[1,R^2]+[r,R^2]\times\{2+4R+R^3+2R^4\}
\]
herleiten. Derselbe Ausdruck ergibt sich aus Formel III aus Artikel~11, nämlich
\[
\begin{array}{l@{\quad}r@{\pdot}r@{\pdot}r@{\pdot}r@{\pdot}r@{\pdot}r@{\pdot}r@{\pdot}r@{\pdot}r}
\text{Werte von }z\ \ldots&2&3&4&5&6&7&8&9&10\\
\operatorname{ind}z\ \ldots&1&8&2&4&9&7&3&6&5\\
\operatorname{ind}(n+1-z)\ \ldots&5&6&3&7&9&4&2&8&1\\
\multicolumn{10}{l}{\text{kleinste Reste von }\operatorname{ind}z+\operatorname{ind}(n+1-z)}\\
\quad\text{zum Modul }5\ \ldots&1&4&0&1&3&1&0&4&1
\end{array}
\]

Auf ganz dieselbe Weise finden wir
\[
\begin{aligned}
2^0.\qquad[r,R^2]\pdot[r,R]&=[1,R^3]+[r,R^3]\times\{2+R+4R^2+2R^3\}\\
3^0.\qquad[r,R^3]\pdot[r,R]&=[1,R^4]+[r,R^4]\times\{2+4R+R^3+2R^4\}
\end{aligned}
\]
Schließlich wird
\[
4^0.\qquad[r,R^4]\pdot[r,R]=[1,1]+[r,1]\times\{1+2R+2R^2+2R^3+2R^4\}
\]
Multiplizieren wir Gleichung $1^0$ mit $[r,R]$ und setzen für $[r,R^2]\pdot[r,R]$ seinen Wert aus $2^0$ ein und verwenden außerdem
\[
[1,R^2]\pdot[r,R]=[r,R^3]\pdot\{2+2R+2R^2+2R^3+2R^4\}
\]
so erhalten wir
\[
\begin{aligned}
\relax[r,R]^3&=[1,R^3]\times\{2+4R+R^3+2R^4\}\\
&\quad+[r,R^3]\times\{12+22R+18R^2+24R^3+15R^4\}
\end{aligned}
\]
und ebenso
\[
\begin{aligned}
\relax[r,R]^4&=[1,R^4]\times\{12+22R+18R^2+24R^3+15R^4\}\\
&\quad+[r,R^4]\times\{164+170R+205R^2+180R^3+190R^4\}\\
\relax[r,R]^5&=[1,1]\times\{164+170R+205R^2+180R^3+190R^4\}\\
&\quad+[r,1]\times\{1836+1830R+1795R^2+1820R^3+1810R^4\}\\
&=1640+1700R+2050R^2+1800R^3+1900R^4\\
&\quad+(1836+1830R+1795R^2+1820R^3+1810R^4)(s-1)\\
&=918Ss-98S-(6R+41R^2+16R^3+26R^4)s\\
&\quad+66R+451R^2+176R^3+286R^4
\end{aligned}
\]

\SourceLeaf{266}
\Article{14}
Die Rechnung in den vorhergehenden Artikeln\textsuperscript{'} ist so durchgeführt, dass sie sich auf sämtliche Werte von $r$ und $R$ erstrecken lässt. Sie wird erheblich kürzer, wenn wir $R$ von Anfang an als eigentliche Wurzel der Gleichung $x^\beta-1=0$ betrachten. Unter dieser Voraussetzung lässt sich das Produkt $[r,R^\mu]\times[r,R^\nu]$ auf die Form $B[r,R^{\mu+\nu}]$ bringen, sofern $\mu+\nu$ nicht durch $\beta$ teilbar ist; ist dagegen $\mu+\nu$ durch $\beta$ teilbar, so wird jenes Produkt $=(n-1)R^{\frac12(n-1)\mu}+[r,1]\Sigma R^{\mu\operatorname{ind}x}$, wenn wir für $\operatorname{ind}x$ alle Zahlen $0,1,2,3\dotsfour n-2$ mit Ausnahme von $\frac12(n-1)$ einsetzen. Daraus folgt leicht, dass in diesem Fall, sofern $\mu$ und damit auch $\nu$ nicht durch $\beta$ teilbar sind,
\[
[r,R^\mu]\pdot[r,R^\nu]=R^{\frac12(n-1)\mu}\{n-1-[r,1]\}
\]
gilt, also $=0$ für $r=1$ und $=nR^{\frac12(n-1)\mu}$ für jeden anderen Wert von $r$. Da ferner $R^{\frac12(n-1)\mu}$ den Wert $=+1$ oder $=-1$ annimmt, je nachdem, ob $\frac{n-1}\beta\pdot\mu$ gerade oder ungerade ist, wird unser Produkt im ersten Fall $=n$, im zweiten $=-n$.

Daraus folgt weiter, dass wir
\[
\begin{aligned}
\relax[r,R]^2&=A'[r,R^2]\\
\relax[r,R^2]\pdot[r,R]&=A''[r,R^3]\\
\relax[r,R^3]\pdot[r,R]&=A'''[r,R^4]
\end{aligned}
\]
setzen können, und so fort bis
\[
[r,R^{\beta-2}]\pdot[r,R]=A^{(\beta-2)}[r,R^{\beta-1}]
\]
woraus sich
\[
\begin{aligned}
\relax[r,R]^2&=A'[r,R^2]\\
\relax[r,R]^3&=A'A''[r,R^3]\\
\relax[r,R]^4&=A'A''A'''[r,R^4]
\end{aligned}
\]
ergibt, und so fort. Schließlich gilt
\[
[r,R]^\beta=\pm nA'A''A'''\dots A^{(\beta-2)}
\]
Dabei ist das obere oder das untere Vorzeichen zu wählen, je nachdem, ob $\frac{n-1}\beta$ gerade oder ungerade ist.

Ist der Wert von $[r,R]$ gefunden, so können die übrigen Funktionen
\[
[r,R^2]=\frac{[r,R]^2}{A'},\quad[r,R^3]=\frac{[r,R]^3}{A'A''}\quad\text{usw.}
\]
hier also offenbar viel leichter bestimmt werden als in den Fällen, in denen die Bedingung $\alpha$ $=1$ nicht erfüllt ist,
\SourceLeaf{267}
wie wir bereits in den \emph{Disq. Ar.} (Artikel~360, \textsc{iii}) angemerkt haben. Eine eingehendere Betrachtung der Natur der Funktionen $A',A''$ usw. erleichtert diese Operationen noch mehr.

\Article{15}
In Artikel~9 haben wir gezeigt, dass sich der Wert der Funktion $[r,R]$ auf die Form $\surd n(\cos f+i\sin f)$ bringen lässt; ebenso lassen sich die Funktionen $[r,R^2]$, $[r,R^3]$ usw. bis $[r,R^{\beta-1}]$ auf eine ähnliche Form bringen. Setzen wir
\[
\begin{aligned}
\relax[r,R]&=\surd n(\cos f'+i\sin f')\\
\relax[r,R^2]&=\surd n(\cos f''+i\sin f'')\\
\relax[r,R^3]&=\surd n(\cos f'''+i\sin f''')\\
&\text{usw.}
\end{aligned}
\]
so wird
\[
\begin{aligned}
A'&=\surd n(\cos(2f'-f'')+i\sin(2f'-f''))\\
A''&=\surd n(\cos(f'+f''-f''')+i\sin(f'+f''-f'''))\\
A'''&=\surd n(\cos(f'+f'''-f'''')+i\sin(f'+f'''-f''''))\\
&\text{usw.}
\end{aligned}
\]
Daraus ist ersichtlich: Werden die Funktionen $A',A'',A'''$ usw. auf die Formen
\[
\begin{aligned}
A'&=a'(\cos b'+i\sin b')\\
A''&=a''(\cos b''+i\sin b'')\\
A'''&=a'''(\cos b'''+i\sin b''')\\
&\text{usw.}
\end{aligned}
\]
gebracht, und zwar so, dass alle $a',a'',a'''$ usw. positiv sind, so gilt
\[
\begin{aligned}
a'&=a''=a'''\text{ usw.}=\surd n\\
f'&=\frac1\beta(b'+b''+b'''+\text{usw.}+b^{(\beta-2)})
\end{aligned}
\]
wenn $\frac{n-1}\beta$ gerade ist, oder
\[
f'=\frac1\beta(180^0+b'+b''+\text{usw.}+b^{(\beta-2)})
\]
wenn $\frac{n-1}\beta$ ungerade ist, und dann
\[
\begin{aligned}
\relax[r,R]&=\surd n(\cos f'+i\sin f')\\
\relax[r,R^2]&=\surd n(\cos(2f'-b')+i\sin(2f'-b'))\\
\relax[r,R^3]&=\surd n(\cos(3f'-b'-b'')+i\sin(3f'-b'-b''))\\
&\text{usw.}
\end{aligned}
\]
\SourceLeaf{268}
Schließlich gilt nach den Formeln aus Artikel~8
\[
\begin{aligned}
\left(\frac{n-1}\beta,1\right)&=-\frac1\beta+\frac{\surd n}\beta\{\cos f''+\cos(2f'-b')+\cos(3f'-b'-b'')+\text{usw.}\\
&\qquad+\cos((\beta-1)f'-b'-b''-b'''-\text{usw.}-b^{(\beta-2)})\}\\
&\quad+\frac{i\surd n}\beta\{\sin f'+\sin(2f'-b')+\sin(3f'-b'-b'')+\text{usw.}\}
\end{aligned}
\]
Ebenso ergeben sich die Werte der Funktionen $\left(\frac{n-1}\beta,g\right)$, $\left(\frac{n-1}\beta,g^2\right)$, $\left(\frac{n-1}\beta,g^3\right)$ usw., wenn wir in dieser Formel für $f'$ der Reihe nach $f'-\frac{360^0k}\beta$, $f'-2\frac{360^0k}\beta$, $f'-3\frac{360^0k}\beta$ usw. einsetzen und dabei $R=\cos\frac{360^0k}\beta+i\sin\frac{360^0k}\beta$ voraussetzen.

\clearpage
\Article{16}
Eine weitere Vereinfachung ergibt sich aus der folgenden Beobachtung. Da nach Artikel~14
\[
\pm[r,R][r,R^{\beta-1}]=[r,R^2][r,R^{\beta-2}]=\pm[r,R^3][r,R^{\beta-3}]\text{ usw.}=n
\]
gilt, wobei [im ersten, dritten usw. Produkt] das obere oder untere Vorzeichen zu wählen ist, je nachdem, ob $\frac{n-1}\beta$ gerade oder ungerade ist, muss im ersten Fall
\[
\cos(f'+f^{(\beta-1)})=\cos(f''+f^{(\beta-2)})=\cos(f'''+f^{(\beta-3)})\text{ usw.}=1
\]
im zweiten
\[
-\cos(f'+f^{(\beta-1)})=\cos(f''+f^{(\beta-2)})=-\cos(f'''+f^{(\beta-3)})\text{ usw.}=1
\]
und in beiden Fällen
\[
\sin(f'+f^{(\beta-1)})=\sin(f''+f^{(\beta-2)})=\sin(f'''+f^{(\beta-3)})\text{ usw.}=0
\]
gelten. Daher dürfen wir im ersten Fall
\[
f^{(\beta-1)}=-f',\quad f^{(\beta-2)}=-f'',\quad f^{(\beta-3)}=-f'''\quad\text{usw.}
\]
im zweiten
\[
f^{(\beta-1)}=180^0-f',\quad f^{(\beta-2)}=-f'',\quad f^{(\beta-3)}=180^0-f'''\quad\text{usw.}
\]
setzen. Daraus folgt aber, dass im ersten Fall
\[
\begin{aligned}
b^{(\beta-2)}&=b',&b^{(\beta-3)}&=b'',&b^{(\beta-4)}&=b'''\quad\text{usw.}\\
A^{(\beta-2)}&=A',&A^{(\beta-3)}&=A'',&A^{(\beta-4)}&=A'''\quad\text{usw.}
\end{aligned}
\]
\Needspace{21\baselineskip}
im zweiten dagegen
\SourceLeaf{269}
\[
\begin{aligned}
b^{(\beta-2)}&=b'-180^0,&b^{(\beta-3)}&=b''+180^0,&b^{(\beta-4)}&=b'''-180^0\quad\text{usw.}\\
A^{(\beta-2)}&=-A',&A^{(\beta-3)}&=-A'',&A^{(\beta-4)}&=-A'''\quad\text{usw.}
\end{aligned}
\]
gilt, sodass sich die Anzahl der Funktionen $A',A'',A'''$ usw. auf die Hälfte verringert. Daraus ergibt sich weiter im ersten Fall
\begingroup\postdisplaypenalty=10000
\[
f'=\frac1\beta(2b'+2b''+\text{usw.}+2b^{(\frac12\beta-1)})
\]
\endgroup
\begingroup\postdisplaypenalty=10000
\[
\begin{aligned}
\left(\frac{n-1}\beta,1\right)&=-\frac1\beta+\frac{\surd n}\beta\{2\cos f'+2\cos(2f'-b')\\
&\qquad+2\cos(3f'-b'-b'')+\text{usw.}\\
&\qquad+2\cos((\tfrac12\beta-1)f'-b'-b''-\text{usw.}-b^{(\frac18\beta-2)})\\
&\qquad+\cos(\tfrac12\beta f'-b'-b''-\text{usw.}-b^{(\frac12\beta-1)})\}
\end{aligned}
\]
\endgroup
(wobei das letzte Glied offenbar $=\cos0=1$ ist) oder
\begingroup\postdisplaypenalty=10000
\[
f'=\frac1\beta(2b'+2b''+\text{usw.}+2b^{(\frac12(\beta-3))}+b^{(\frac12(\beta-1))})
\]
\endgroup
\[
\begin{aligned}
\left(\frac{n-1}\beta,1\right)&=-\frac1\beta+\frac{\surd n}\beta\{2\cos f'+2\cos(2f'-b')\\
&\qquad+2\cos(3f'-b'-b'')+\text{usw.}\\
&\qquad+2\cos(\tfrac12(\beta-1)f'-b'-b''-\text{usw.}-b^{(\frac12(\beta-3))})\}
\end{aligned}
\]
je nachdem, ob $\beta$ gerade oder ungerade ist; und im zweiten Fall

\begingroup\postdisplaypenalty=10000
\[
f'=\frac1\beta(2b'+2b''+\text{usw.}+2b^{(\frac12\beta-1)})
\]
\endgroup
\[
\begin{aligned}
\left(\frac{n-1}\beta,1\right)&=-\frac1\beta+\frac{\surd n}\beta\{2\cos(2f'-b')+2\cos(4f'-b'-b''-b''')+\text{usw.}\\
&\qquad+2\cos((\tfrac12\beta-2)f'-b'-b''-\text{usw.}-b^{(\frac12\beta-3)})\\
&\qquad+\cos(\tfrac12\beta f'-b'-b''-\text{usw.}-b^{(\frac12\beta-1)})\}\\
&\quad+i\frac{\surd n}\beta\{2\sin f'+2\sin(3f'-b'-b'')+\text{usw.}\\
&\qquad+2\sin((\tfrac12\beta-1)f'-b'-b''-\text{usw.}-b^{(\frac12\beta-2)})\}
\end{aligned}
\]
\Needspace{32\baselineskip}
oder
\[
f'=\frac1\beta(2b'+2b''+\text{usw.}+2b^{(\frac12\beta-1)}+180^0)
\]
\[
\begin{aligned}
\left(\frac{n-1}\beta,1\right)&=-\frac1\beta+\frac{\surd n}\beta\{2\cos(2f'-b')+2\cos(4f'-b'-b''-b''')+\text{usw.}\\
&\qquad+2\cos((\tfrac12\beta-1)f'-b'-b''-\text{usw.}-b^{(\frac12\beta-2)})\}\\
&\quad+i\frac{\surd n}\beta\{2\sin f'+2\sin(3f'-b'-b'')+\text{usw.}\\
&\qquad+2\sin((\tfrac12\beta-2)f'-b'-b''-\text{usw.}-b^{(\frac12\beta-3)})\\
&\qquad+\sin(\tfrac12\beta f'-b'-b''-\text{usw.}-b^{(\frac12\beta-1)})\}
\end{aligned}
\]
je nachdem, ob $\frac12\beta$ gerade oder ungerade ist. Für die übrigen Perioden mit $\frac{n-1}\beta$ Gliedern gilt dasselbe, was wir oben angemerkt haben. Allgemein folgt daraus also, dass zur Bestimmung dieser Perioden die Teilung des ganzen Kreises in $\beta$ Teile erforderlich ist, von der die
\SourceLeaf{270}
Konstruktion der Winkel $b',b'',b'''$ usw. rational abhängt; ferner die Teilung des Winkels $b'+b''+b'''+\text{usw.}$ in $\beta$ Teile; schließlich die Quadratwurzel $\surd n$. Setzen wir sogleich $\beta=\frac12(n-1)$, so stimmen jene Perioden offenbar mit den doppelten Kosinus der Winkel $\frac{360^0}n$, $2\frac{360^0}n$, $3\frac{360^0}n$ usw. bis $\frac12(n-1)\frac{360^0}n$ überein. Daher hängt die Teilung des Kreises in $n$ Teile von der Teilung des ganzen Kreises in $\frac12(n-1)$ Teile ab, von der Teilung eines nach Ausführung jener Teilung konstruierbaren Winkels in $\frac12(n-1)$ Teile und von der Wurzelgröße $\surd n$. Wollen wir bis zu den Sinus der Winkel $\frac{360^0}n$ usw. fortgehen, so ist eine weitere Operation nötig.
